\section{Síntesis y ampliación}

Esta sección amplía el frontier scaling, que deja de ser «cuanto más grande el modelo, mejor» y se convierte en un operating system de research, systems, safety y product:

\begin{enumerate}
  \item \textbf{La scaling law es un contrato de predicción}: primero se hacen el fit y el residual con runs a varias escalas, y solo después se aprueba el frontier commitment;
  \item \textbf{El compute multiplier requiere atribución}: los cambios de model, data, optimization y systems deben validarse con una misma base de costo;
  \item \textbf{La rare failure debe incorporarse al diseño}: el failure domain, el checkpoint, el replay y el recovery tax determinan el useful progress;
  \item \textbf{El entrenamiento requiere observabilidad en cuatro planos}: la telemetry de model, optimizer, data y system explica en conjunto cada anomaly;
  \item \textbf{La guardia global requiere decision rights}: el handoff packet, el owner y los stop criteria importan más que «haya alguien conectado»;
  \item \textbf{El post-training optimiza un proxy}: RLHF/CAI/RLAIF amplían la supervisión del comportamiento, pero todos requieren evaluation y gobernanza independientes;
  \item \textbf{La evaluation depende de la elicitation}: task, tool, effort, environment, judge e uncertainty determinan qué significa una puntuación;
  \item \textbf{Las safeguards escalan con la capability}: las versiones de la RSP evolucionan, pero el principio de evaluar primero, cumplir los required controls y después entrenar/desplegar no cambia;
  \item \textbf{La interpretabilidad es evidencia complementaria}: una feature/circuit requiere causal validation y no puede demostrar la seguridad por sí sola;
  \item \textbf{La API es un contract externo}: en el chat se puede experimentar con rapidez; la API debe tener version, migrate, deprecate y support.
\end{enumerate}

\begin{importantbox}{Tarea: Frontier Model Program}
Diseña un plan de entrenamiento y lanzamiento de 8 semanas: presenta tres niveles de pilot scale, con su loss fit y su uncertainty; enumera el failure budget de workers/network/storage/scheduler; define el checkpoint/replay y el handoff follow-the-sun; diagrama el preentrenamiento, RLHF/RLAIF, la capability/safety evaluation, el safeguard gate, el chat canary y el API release; por último, presenta la model deprecation, el shadow traffic y el rollback. Cada gate debe tener owner, evidence y stop condition.
\end{importantbox}

\subsection{Lecturas adicionales}

{\small
\begin{itemize}[nosep,leftmargin=*]
  \item \href{https://arxiv.org/abs/2001.08361}{\textbf{Scaling Laws for Neural Language Models}}: escalamiento empírico y métodos de predicción.
  \item \href{https://arxiv.org/abs/2203.02155}{\textbf{InstructGPT}}: el pipeline de RLHF y sus limitaciones.
  \item \href{https://www.anthropic.com/research/constitutional-ai-harmlessness-from-ai-feedback}{\textbf{Constitutional AI}}: principles, critique/revision y RLAIF.
  \item \href{https://www.anthropic.com/research/evaluating-ai-systems}{\textbf{Challenges in evaluating AI systems}}: diseño de evaluaciones y su validez en el mundo real.
  \item \href{https://www.anthropic.com/responsible-scaling-policy/rsp-v3-0}{\textbf{Anthropic RSP v3.0}}: el capability-threshold framework vigente en 2026.
  \item \href{https://www.anthropic.com/research/towards-monosemanticity-decomposing-language-models-with-dictionary-learning}{\textbf{Towards Monosemanticity}}: dictionary learning y features.
  \item \href{https://docs.anthropic.com/en/docs/about-claude/model-deprecations}{\textbf{Claude model lifecycle}}: active, legacy, deprecated y retired.
\end{itemize}
}

\end{document}
